\section{Complex Tori and the Jacobian}
\label{sec:complex-tori-jacobians}

Let \(V\) be a finite-dimensional complex vector space and \(\Lambda\subset V\)
a full lattice. Then \(V/\Lambda\) is a compact complex Lie additive group.
The quotient map is a local biholomorphism after restricting to sufficiently
small balls avoiding nonzero lattice translates.

\begin{proposition}[Complex torus package]
\label{prop:complex-torus-package}
The quotient \(V/\Lambda\) has an additive commutative group structure,
Hausdorff topology, compactness, complex charted-space structure, complex
manifold structure, and compatible Lie additive group structure.
\lean{JacobianChallenge.ComplexTorus.FullComplexLattice, JacobianChallenge.ComplexTorus.quotient, JacobianChallenge.ComplexTorus.quotient_addCommGroup, JacobianChallenge.ComplexTorus.quotient_topologicalSpace, JacobianChallenge.ComplexTorus.quotient_t2Space, JacobianChallenge.ComplexTorus.quotient_compactSpace}
\leanok
\end{proposition}

\begin{layman}
Take a finite-dimensional complex vector space \(V\) — a few copies of
the complex plane — and a \emph{full lattice} \(\Lambda\) inside it: a
discrete grid of points whose integer combinations spread out evenly in
every direction, like the vertices of a tiled wallpaper pattern in
real-dimension \(2g\). Glue the space to itself by declaring two points
equivalent whenever they differ by a lattice vector. The result is a
\emph{complex torus}: in dimension one, just the surface of an ordinary
donut; in higher dimension, a higher-genus generalization that's still
compact, smooth, and inherits an addition law from \(V\). The
proposition packages all the structural facts at once: addition
commutes, the topology is Hausdorff and compact, the smooth structure
makes it a complex manifold, and addition itself is smooth. (You can
think of the building blocks like a Lego set: snap together vector
space, group, topology, manifold, and Lie group.) Why bother? When the
Jacobian is built later as ``dual of holomorphic 1-forms modulo period
lattice,'' we just plug \(V = \HdualX\) and \(\Lambda = \) period
lattice into this proposition and inherit every structural property
the challenge's anti-hack theorems demand.
\end{layman}

\begin{proof}
The group and topology are induced by the quotient of the additive topological
group \(V\). Closedness and discreteness of \(\Lambda\) give Hausdorffness and
local injectivity of the quotient map. A compact fundamental parallelepiped
covers \(V/\Lambda\), proving compactness. Local charts are obtained by
choosing a ball \(B(v,r)\) on which the quotient map is injective; transition
maps are translations in \(V\), hence holomorphic and smooth. Addition and
negation are induced from smooth maps on \(V\).
\end{proof}

\begin{proof}[Construction of Definition~\ref{def:analytic-jacobian}]
Apply Proposition~\ref{prop:complex-torus-package} to
\[
  V=\HdualX,\qquad \Lambda=\Pi_X(H_1(X,\Z)).
\]
The full-lattice condition is exactly Theorem~\ref{thm:period-lattice}. This
produces the Jacobian \(\Jac(X)\) with all structural instances required by
the challenge.
\end{proof}

\begin{formalizationnote}
Most complex torus infrastructure is already present in
\code{Jacobian/ComplexTorus}. The challenge-facing carrier uses a
basis-aligned model \(\operatorname{Fin}(g)\to\C\) and a universe bridge via
\code{ULift}. The informal proof treats this as transport along a chosen
linear equivalence \(\HdualX\cong\C^g\).
\end{formalizationnote}
